\documentclass[11pt]{article}
\usepackage[a4paper,margin=1.9cm]{geometry}
\usepackage[T1]{fontenc}
\usepackage{textcomp}
\usepackage[spanish]{babel}
\usepackage{amsmath,amssymb}
\usepackage{lmodern}
\renewcommand{\familydefault}{\sfdefault}
\usepackage{enumitem}

\newcommand{\vect}[2]{\begin{pmatrix}#1\\#2\end{pmatrix}}
\newcommand{\vectiii}[3]{\begin{pmatrix}#1\\#2\\#3\end{pmatrix}}
\usepackage{tikz}
\usetikzlibrary{calc,arrows.meta,patterns,decorations.pathmorphing,3d}
\usepackage[table]{xcolor}
\usepackage{multicol}
\usepackage{parskip}

\definecolor{unalblue}{RGB}{31,73,124}

\newlist{opciones}{enumerate}{1}
\setlist[opciones]{label=(\alph*),itemsep=6pt,topsep=6pt,leftmargin=1.8em,font=\normalfont}
\setlist[enumerate,1]{itemsep=1.4em,topsep=1em}

\newcommand{\examheader}{%
  \noindent
  \begin{minipage}[t]{0.6\linewidth}
    {\small\bfseries Geometría Vectorial y Analítica --- Parcial 3}\\
    {\small Semestre 2015--02}
  \end{minipage}%
  \hfill
  \begin{minipage}[t]{0.35\linewidth}
    \raggedleft\small Escuela de Matemáticas. Universidad Nacional de Colombia, Sede Medellín.
  \end{minipage}\\[2pt]
  \rule{\linewidth}{0.6pt}
}
\newcommand{\examfooter}{%
  \vfill
  \rule{\linewidth}{0.4pt}\\[1pt]
  {\scriptsize\textit{Transcripción digital --- MathUNAL}\hfill\textcolor{unalblue}{\textbf{\thepage}}}
}
\pagestyle{empty}

\newcommand{\nota}[1]{{\small\itshape\color{unalblue!70!black} #1}}

\begin{document}
\examheader

\begin{center}
{\large Tercer Examen Parcial de Geometría}\\[4pt]
{\small 23 de noviembre del 2015}
\end{center}

\vspace{6pt}
\textbf{Puntaje. Sólo para uso Oficial}

\begin{center}
\begin{tabular}{|c|c|c|c|c|c|}
\hline
1 (/30) & 2 (/20) & 3-5 (/30) & 6 (/20) & Total & Nota \\
\hline
 & & & & & \\
\hline
\end{tabular}
\end{center}

\vspace{10pt}
\textbf{Instrucciones:} Antes de empezar verifique que su examen esté completo y que sus datos sean correctos. Por favor verifique que su celular esté apagado. No está permitido el uso de calculadora ni de otros aparatos electrónicos. Solucione las preguntas en los espacios en blanco asignados a cada una de ellas. No está permitido sacar hojas en blanco ni ningún tipo de apuntes durante el examen. \textbf{Duración: 1 hora y 50 minutos.}

\vspace{6pt}
\nota{En cada uno de los literales de los puntos 1 y 2, debe presentar detalladamente el procedimiento para llegar a la solución. NOTA: Respuestas sin procedimiento no se tendrán en cuenta.}

\begin{enumerate}
\item Considere los puntos $X=\vectiii{-1}{0}{0}$, $Y=\vectiii{0}{-1}{0}$ y $Z=\vectiii{0}{0}{-1}$.

\begin{enumerate}[label=(\alph*)]
\item{}[10pt] Encuentre la ecuación general del plano $\mathcal{P}$, que pasa por los puntos $X$, $Y$ y $Z$.
\item{}[10pt] Encuentre las ecuaciones escalares paramétricas de la recta $\mathcal{L}$, que pasa por $O=\vectiii{0}{0}{0}$ y es perpendicular al plano $\mathcal{P}$.
\item{}[10pt] Determine cuál es el punto del plano $\mathcal{P}$ que está más cercano al origen.
\end{enumerate}

\item Sea $\mathcal{L}_1$ la recta que pasa por los puntos $A=\vectiii{0}{2}{0}$ y $B=\vectiii{0}{0}{2}$.

\begin{enumerate}[label=(\alph*)]
\item{}[10pt] Encuentre la ecuación general del plano $\mathcal{P}$, que contiene a la recta $\mathcal{L}_1$ y pasa por el punto $C=\vectiii{-1}{-1}{-1}$.
\item{}[10pt] Se sabe que la recta $\mathcal{L}_2: x-1=\dfrac{y-1}{3}=z-1$, es paralela al plano $\mathcal{P}$. Encuentre la distancia entre $\mathcal{L}_1$ y $\mathcal{L}_2$.
\end{enumerate}
\end{enumerate}

\vspace{6pt}
\nota{En las preguntas 3 a 5 complete los espacios en blanco o elija la (las) opción (opciones) correcta (correctas) según el caso. Marque con una X la letra de su elección. NOTA: En esta sección se califica sólo la respuesta y no se tiene en cuenta el procedimiento.}

\begin{enumerate}
\setcounter{enumi}{2}
\item{}[12pt] En la figura se muestra una parábola con foco $F$, vértice en la intersección de los ejes $x'$, $y'$ y cuya directriz es la recta $\mathcal{L}$. Se sabe que $P$ pertenece a la parábola, que $\operatorname{dist}(P,\mathcal{L})=7$ y que $\|F-Q\|=6$.

\begin{center}
\begin{tikzpicture}[scale=0.9]
\draw[dashed] (-2.6,-1.3) -- (2.6,1.3);
\node at (2.9,1.45) {$\mathcal{L}$};
\draw[-{Latex}] (-2,2) -- (2.2,-2.2) node[right]{$x'$};
\draw[-{Latex}] (-1.2,-2) -- (1.2,2) node[above]{$y'$};
\draw[thick,domain=-1.3:1.5,smooth,variable=\t] plot ({-0.35*\t*\t+0.3},{\t});
\node[left] at (0.3,0) {$F$};
\fill (0.3,0) circle (1.2pt);
\fill (-1.1,-0.55) circle (1.2pt);
\node[below] at (-1.1,-0.6) {$Q$};
\fill (0.05,1.15) circle (1.2pt);
\node[above right] at (0.05,1.15) {$P$};
\end{tikzpicture}
\end{center}

\begin{enumerate}[label=(\alph*)]
\item $\|F-P\|-2\operatorname{dist}(P,\mathcal{L})=$ \underline{\hspace{2.5cm}}.
\item La ecuación de la parábola en el sistema $x',y'$ es: \underline{\hspace{4cm}}
\end{enumerate}

\item{}[9pt] El sistema $(V,U,X)$ de la figura está orientado con la mano derecha y además los vectores son ortogonales entre sí. Se sabe también que $X=-Y$ y que $\|X\|=$ Área Paralelogramo $OVRU$.

\begin{center}
\begin{tikzpicture}[scale=0.9]
\coordinate (O) at (0,0);
\coordinate (V) at (1.6,-0.8);
\coordinate (U) at (2.6,0.3);
\coordinate (R) at ($(V)+(U)$);
\draw[fill=black!15] (O) -- (V) -- (R) -- (U) -- cycle;
\draw[-{Latex},thick] (O) -- (V) node[below right]{$V$};
\draw[-{Latex},thick] (O) -- (U) node[right]{$U$};
\draw[-{Latex},thick] (O) -- (0,2.2) node[above]{$X$};
\draw[-{Latex},thick] (O) -- (0,-1.8) node[below]{$Y$};
\node[left] at (O) {$O$};
\end{tikzpicture}
\end{center}

\begin{enumerate}[label=(\alph*)]
\item Si $\|V\|=1$ entonces $V\times Y=$ \underline{\hspace{2.5cm}}.
\item Si $\|X\|=6$ y $\|U\|=2$ entonces $\|V\|=$ \underline{\hspace{2.5cm}}.
\item Si $\|Y\|=5$ entonces $U\cdot(X\times V)=$ \underline{\hspace{2.5cm}}.
\end{enumerate}

\item{}[9pt] Las rectas $\mathcal{L}_1$, $\mathcal{L}_2$, $\mathcal{L}_3$ de la figura forman un triángulo en el espacio, cuyos directores son los vectores unitarios $D_1$, $D_2$ y $D_3$ respectivamente. Se sabe que $D_1\cdot D_3=-\dfrac{\sqrt{3}}{2}$ y que el ángulo entre $\mathcal{L}_1$ y $\mathcal{L}_2$ es $\dfrac{\pi}{4}$.

\begin{center}
\begin{tikzpicture}[scale=0.9]
\coordinate (P1) at (-1.3,-1.1);
\coordinate (P2) at (1.6,-0.5);
\coordinate (P3) at (-0.5,1.5);
\draw[thick] (P1) -- (P2) -- (P3) -- cycle;
\draw[dotted] ($(P3)!-0.3!(P1)$) -- ($(P1)!-0.3!(P3)$);
\draw[-{Latex},thick] (P1) -- ($(P1)+(-0.9,-0.9)$) node[below]{$D_3$};
\draw[-{Latex},thick] (P3) -- ($(P3)+(-0.7,0.9)$) node[above]{$D_2$};
\draw[-{Latex},thick] (P2) -- ($(P2)+(1.1,0.2)$) node[right]{$D_1$};
\node[below left,font=\small] at (P1) {$\beta$};
\node[above,font=\small] at (P3) {$\alpha$};
\node[right,font=\small] at (P2) {$\gamma$};
\node[below] at (P1) {$\mathcal{L}_1$};
\node[above right] at (P3) {$\mathcal{L}_3$};
\node[right] at (P2) {$\mathcal{L}_2$};
\end{tikzpicture}
\end{center}

\begin{enumerate}[label=(\alph*)]
\item El ángulo entre $\mathcal{L}_1$ y $\mathcal{L}_3$ es \underline{\hspace{2.5cm}}.
\item $D_1\cdot D_2=$ \underline{\hspace{2.5cm}}.
\item $\alpha=$ \underline{\hspace{2.5cm}}.
\end{enumerate}

\item Una elipse $\mathcal{E}$ está dada por la ecuación
\begin{equation}
\frac{7}{16}x^2-\frac{6\sqrt{3}}{16}xy+\frac{13}{16}y^2-1=0.
\end{equation}

\begin{enumerate}[label=(\alph*)]
\item{}[7pt] Sabiendo que $M=\begin{pmatrix}\frac{7}{16} & -\frac{3\sqrt{3}}{16}\\ -\frac{3\sqrt{3}}{16} & \frac{13}{16}\end{pmatrix}$ se puede factorizar en la forma $QDQ^T$ como
$$M=\begin{pmatrix}\frac{\sqrt{3}}{2} & -\frac{1}{2}\\ \frac{1}{2} & \frac{\sqrt{3}}{2}\end{pmatrix}\begin{pmatrix}\frac{1}{4} & 0\\ 0 & 1\end{pmatrix}\begin{pmatrix}\frac{\sqrt{3}}{2} & \frac{1}{2}\\ -\frac{1}{2} & \frac{\sqrt{3}}{2}\end{pmatrix},$$
escriba la ecuación (1) en la forma $\widehat{A}(x')^2+\widehat{C}(y')^2+F=0$ indicando explícitamente el valor de los coeficientes. Justifique su respuesta.
\item{}[7pt] Halle los vértices de la elipse $\mathcal{E}$ en las coordenadas $x'y'$. Justifique su respuesta.
\item{}[6pt] Si se sabe que $P=\vect{p}{q}$ pertenece a la elipse $\mathcal{E}$, ¿cuál es la suma de las distancias de $P$ a los focos de $\mathcal{E}$? Justifique su respuesta.
\end{enumerate}

\end{enumerate}

\examfooter
\end{document}
